\section{Philosophy of Language}

\section{Course Description}

According to a traditional philosophical picture, words express ideas and refer to things in the world, and sentences are true or false depending on the facts about those things. In the twentieth century, philosophers began to question and criticize this picture of meaning, by asking such
questions as: Are there different kinds of meaning? What is done by means of language, besides referring
to things in the world and stating facts? How do speakers’ intentions influence the meaning of their words?
We will examine some of the more sophisticated views of language that emerged, focusing especially on
the idea that in some sense, meaning lies in how language is used. But this view faces its own challenges.
Which aspects of use are relevant to meaning? How do we know when words are used correctly? How
do we know what others use words to refer to? These challenges will lead us to re-examine the traditional
picture. Does the idea that we use language to speak about things in the world make any sense at all? If so,
how can we make sense of it?


\section{Course Objectives}

Students will be able to:
1. Explain important concepts in philosophy of language such as “truth-condition”,
“proposition”, “rigid designator”, “speech act”, “common ground”, “conversational
implicature”, “conventional implicature”, “presupposition”, etc.
2. Understand important classic papers in the philosophy of language or related fields.
3. Charitably reconstruct arguments from texts and evaluate them.
4. Apply conceptual and logical tools in philosophy of language to practical problems.

Week 1 Overview & Basic Concepts in the Philosophy of Language

  The Origins of Truth-Conditional Semantics

  Week 1
  Gottlob Frege, “On Sense and Nominatum”, (pp. 29-35  Donde termina el párrafo con ...referencia)

    Notar para el glosario: relación, igualdad, conocimiento efectivo, sentido, referencia, nombre propio, sentido indirecto, representación

Week 2

Gottlob Frege, “On Sense and Nominatum”, (pp. 35-49.)

Esto surge de una imperfección del lenguaje, de la que tampoco está com-
pletamente libre por lo demás el lenguaje simbólico del análisis; también aquí
pueden aparecer combinaciones de signos que producen la apariencia de que
se refieren a algo pero que, por lo menos hasta ahora, carecen de referencia,
como, por ejemplo, las series divergentes infinitas. Esto puede evitarse por
medio de una estipulación especial al efecto de que, por ejemplo, las series
divergentes infinitas deban referirse al número 0. De un lenguaje lógicamente
perfecto (conceptografía) ha de reclamarse que cada expresión que se ha for-
mado como nombre propio de modo gramaticalmente correcto, a partir de
signos ya introducidos, designe también de hecho un objeto, y que no se intro-
duzca ningún nuevo signo como nombre propio, sin que tenga asegurada una
referencia. En los textos de lógica se advierte sobre la ambigiedad de las expre-
siones como fuente de errores lógicos. Como mínimo, tengo por igualmen-


En el caso de estas oraciones, son posibles, por lo demás, concepciones ligeramente
distintas. El sentido de la oración «Después de que Schleswig-Holstein se hubo separado de
Dinamarca, Prusia y Austria se enemistaron», puede reproducirse de la forma: «Después de
la separación de Schleswig-Holstein de Dinamarca, Prusia y Austria se enemistarom». De acuer-
do con esta concepción, resulta en verdad suficientemente claro que el pensamiento de que
Schleswig-Holstein se separó alguna vez de Dinamarca no ha de tomarse como parte de este
sentido, sino que ese pensamiento es la presuposición necesaria para que la expresión «des-
pués de la separación de Schleswig-Holstein de Dinamarca» tenga alguna referencia. Desde
luego, nuestra oración puede concebirse de tal manera que se esté diciendo con ella que Schles-
wig-Holstein se separó alguna vez de Dinamarca. Tenemos entonces un caso que ha de con-
siderarse más tarde. Para darse cuenta más claramente de la diferencia, pongámonos en el
pellejo de un chino que, debido a su escaso conocimiento de la historia europea, tiene por
falso el que Schleswig-Holstein se haya separado alguna vez de Dinamarca. Éste no tendría
a nuestra oración, concebida del primer modo, ni por verdadera ni por falsa, sino que le nega-
ría cualquier referencia, puesto que la oración subordinada carece de ella. Esta última sólo
daría, aparentemente, una determinación temporal. Si, por el contrario, concibiese nuestra
oración de acuerdo con el segundo modo, encontraría expresado en ella un pensamiento que
tendría por falso, junto a una parte que, para él, carecería de referencia.


Cuando ha de indicarse un instante de manera indeterminada en el ante-
cedente y el consecuente de una oración condicional, esto ocurre frecuen-
temente por medio del mero uso del tiempo presente del verbo que, en tal
caso, no conlleva la designación del presente temporal. Esta forma grama-
tical es entonces el componente que indica de manera indeterminada en la
oración principal y en la subordinada. «Cuando el Sol está en el trópico de
Cáncer, tenemos el día más largo en el hemisferio norte» es un ejemplo de
esto. Aquí es también imposible expresar el sentido de la subordinada en
una principal, puesto que este sentido no es ningún pensamiento completo;
pues si dijésemos: «El Sol está en el trópico de Cáncer», estaríamos desig-
nando nuestro presente y, con ello, alteraríamos el sentido. Aún menos es
el sentido de la principal un pensamiento; sólo el todo compuesto de la ora-
ción principal y de la subordinada contiene tal pensamiento. Por lo demás,
también se pueden indicar de manera indeterminada diversos componentes
comunes del antecedente y el consecuente de una oración condicional.


Si encontramos que, por lo general, el valor cognoscitivo de «a = a» y
«a = b» es diferente, esto se explica diciendo que, para el valor cognosciti-
vo, el sentido de la oración, a saber: el pensamiento expresado por ella, no
viene menos al caso que su referencia, esto es: su valor de verdad. Ahora
bien, si a = b, entonces la referencia de «b» es ciertamente la misma que la
de «a» y, por consiguiente, también el valor de verdad de «a = b» es el mismo
que el de «a = a». Con todo, el sentido de «b» puede ser distinto del senti-
do de «a», y con ello también el pensamiento expresado en «a = b» será dife-
rente del expresado por «a = a»; en ese caso, ambas oraciones tampoco tie-
nen el mismo valor cognoscitivo. Si entendemos por «juicio», como lo hemos
hecho más arriba, el avance del pensamiento a su valor de verdad, diremos
también entonces que los juicios son diferentes.


%Se comienza a partir en temas en 34 se vuelca al sentido y refProbablemente, uno se decidiría por esta opción. En el otro caso la situa-
ción estaría bastante enredada: tendríamos entonces más pensamientos sim-
ples que oraciones. Si sustituimos también ahora la oración
«Napoleón cayó en la cuenta del peligro para su flanco derecho»
por otra con el mismo valor de verdad, por ejemplo, por
«Napoleón tenía ya más de cuarenta y cinco años»,
se habría alterado con ello no sólo nuestro primer pensamiento, sino tam-
bién nuestro tercero, y por ello podría cambiar su valor de verdad, es decir:
si su edad no era la razón de la decisión de dirigir su guardia contra el ene-
migo. Á partir de esto, puede percibirse por qué no siempre se pueden reem-
plazar entre sí las oraciones con el mismo valor de verdad. La oración expre-
sa entonces, en virtud de su conexión con otra, más que por sí sola.
Consideremos ahora casos donde tal cosa sucede regularmente. En la
oración
«Bebel se imagina que la devolución de Alsacia-Lorena aplacará los
deseos de venganza de Francia» erencia de oraciones, no de nombres.


En la expresión del primer pensamiento, las palabras de la oración sub-
ordinada tienen su referencia indirecta, mientras que las mismas palabras
tienen su referencia habitual en la expresión del segundo pensamiento. Vemos
a partir de esto que la subordinada de nuestra oración compuesta originaria
ha de tomarse ciertamente como doble con diferentes referencias, de las cua-
les una es un pensamiento, la otra un valor de verdad. Ahora bien, puesto
que el valor de verdad no es la referencia total de la oración subordinada,
no podemos simplemente reemplazar ésta por otra con el mismo valor de

verdad. Sucede algo similar con expresiones como «saber», «darse cuen-
ta», «se sabe».


Con una oración subordinada causal y su oración principal co-
rrespondiente, expresamos diversos pensamientos que, sin embargo, no
corresponden a las oraciones tomadas de modo aislado. En la oración
«Puesto que el hielo es menos denso que el agua, flota en el agua»
tenemos:
1.
El hielo es menos denso que el agua.
2.
Si algo es menos denso que el agua, flota en el agua.
3.
El hielo flota en el agua.
El tercer pensamiento no necesitaba quizás ser introducido explícitamente,
puesto que está contenido en los dos primeros. En cambio, ni combinando
el primero con el tercero ni el segundo con el tercero, se lograría el sentido
de nuestra oración. Puede verse ahora que en nuestra oración subordinada
«puesto que el hielo es más denso que el agua»
se expresa tanto nuestro primer pensamiento como una parte de nuestro segun-
do. De aquí viene el que no podamos reemplazar sin más nuestra subordi-
nada por otra con el mismo valor de verdad; pues con ello se alteraría tam-
bién nuestro segundo pensamiento y podría fácilmente verse afectado
también su valor de verdad.
El asunto es similar en la oración
«Si el hierro fuera menos denso que el agua, flotaría en el agua».
Tenemos aquí los dos pensamientos siguientes: que el hierro no es menos
denso que el agua, y que algo flota en el agua si es menos denso que el agua.



%En 36 es donde aparece el prototipo de la teoría deflacionista de la verdad y al final aparece el criterio \textit{salva veritate}

%37 oraciones subordinadas parte de actitudes proposicionales (contextos opacos)

%39 referencia indirecta, uso de órdenes (actos de habla)

%40 Descripciones definidas.

21 Gottlob Frege, “On Sense and Nominatum”, (pp. 35-)
Frege, “On Sense and Reference”  (pp. 217-229)


Russell, “On Denoting” (pp. 230-238)

week 3
Strawson, “On Referring”  (pp. 246-260)
Donnellan, “Reference and Definite Descriptions” (pp. 265-277)

Week 4
Grice, “Logic and Conversation”
Austin, How to Do Things with Words, Lectures 8-12
Kaplan Indexicals

Week 5
Berstler, “What’s the Good of Language? On the Moral Distinction Between Lying and Misleading”

%Elisabeth Camp, “Contextualism, Metaphor, and What is Said”

%Saul Kripke, Naming and Necessity (selections)

%Week 11 Ruth Barcan Marcus, “Modalities and Intensional %Languages” (selections)

%David Kaplan, “Demonstratives” (selections)


\subsection{Grade Distribution}

The overall grade is determined by the following:
\begin{itemize}
\item Participation 70%
\item Grade 10%
\item Final essay 20%
\end{itemize}

\section{Skills developed}

(i) to gain familiarity with the history of analytic
philosophy as it relates to the philosophy of language (ii)
to gain familiarity with the topics of interest in analytic
philosophy of language (iii) to improve your ability to
philosophically engage with the issues underlying (i) and
(ii) (iv) to acquire the ability to extend the course
material to different cultural and linguistic contexts

Philosophy is an activity that we do, and active
participation in philosophy is the best way to learn to do
philosophy. You are expected to interact with me and with
other students inside and outside of class. It’s important
to note, though, that active participation is more than just
being vocal; it requires carefully thinking through issues
and engaging with your peers, often by listening to,
supporting, clarifying, or justifying their comments. Doing
philosophy is not just about expressing your own ideas, but
is just as much about engaging with the ideas of others.
Metaphorically speaking, the ideal philosophical discussion
is less like a game of ping-pong and more like a soccer
(“football”) match. You will be graded on the extent to
which you follow this model of active participation.

Day 11: Davidson’s “A Nice Derangement of Epitaphs” Day 12:
Austin’s “Performative Utterances


Date Reading Selection(s)
7 Course introduction


Se comienza a partir en temas en 34 se vuelca al sentido y referencia de oraciones, no de nombres.

En 36 es donde aparece el prototipo de la teoría deflacionista de la verdad y al final aparece el criterio \textit{salva veritate}

37 oraciones subordinadas parte de actitudes proposicionales (contextos opacos)

39 referencia indirecta, uso de órdenes (actos de habla)
40 descripciones definidas, obviamente russell es idiota frege apoya strawson



Mon, Feb 9 Bertrand Russell, “Descriptions” (pp. 239-245)
Wed, Feb 11 P.F. Strawson, “On Referring” ;
Keith Donnellan, “Reference and Definite Descriptions”
Mon, Feb 16 Saul Kripke, “Naming and Necessity” (pp. 290-305)
Wed, Feb 18 Hilary Putnam, “Meaning and Reference” (pp. 306-313)

David Kaplan, “Quantifying In” (pp. 399-419)

Jon Barwise and John Perry, “Semantic Innocence and Uncompromising Situations” (pp. 420-432)

David Kaplan, “On the Logic of Demonstratives” (pp. 357-365)

Robert Stalnaker, “Semantics for Belief” (pp. 460-468

H.P. Grice, “Meaning” (pp. 108-113)
Mon, Apr 13 Carl Hempel, “Empiricist Criteria of Cognitive Significance: Problems and Changes” (pp. 50-62)
Wed, Apr 15 W.V. Quine, “Two Dogmas of Empiricism” (pp. 63-76)
Mon, Apr 20 W.V. Quine, “Translation and Meaning” (pp. 546-575)
Wed, Apr 22 Alfred Tarski, “The Semantic Conception of Truth and the Foundations of Semantics” (pp. 85-107)
Mon, Apr 27 Donald Davidson, “Truth and Meaning” (pp. 114-125) and “Belief and the Basis of Meaning” (pp. 576-
584)
Wed, Apr 29 John Searle, “Indeterminacy, Empiricism and the First Person” (pp. 596-610)
Mon, May 4 Saul Kripke, “On Rules and Private Language” (pp. 626-638)
Wed, May 6 Ruth Millikan, “Truth Rules, Hoverflies, and the Kripke-Wittgenstein Paradox” (pp. 639-655


dvice
Writing philosophy is difficult, and doing it well takes time and practice. As this is an upper level philosophy course, I assume you've had some practice already, and you're about to get yet more of it in this course. In addition, you may find it useful to consult some of Jim Pryor's excellent tips on philosophical writing. Obviously, these materials are not designed for our writing assignments in particular, but the advice given is generally sound and applicable.

I can't emphasize enough the importance of starting your writing early. The process of writing -- even if only starting with a half-baked idea -- will help you crystalize your thoughts and get clear on what you do and don't understand. That, in turn, will tell you what you need to do next to refine your ideas. Also, starting early allows you the (necessary) luxury of setting out, reconsidering, revising, and developing lines of thought.

And if there's something in the course that you don't understand, come see me about it. Getting the issue cleared up sooner rather than later means that it won't create other problems for you, and will allow you more time to enjoy the warm glow of understanding. This is why God invented office hours. You've paid your tuition, so don't let material go over your head; come and get the education you deserve.

Segment 2 (weeks 3-4): Frege on sense and reference

    Frege, "On Sense and Reference"
    Frege, "The Thought"



	Segment 3 (weeks 5-6): Definite descriptions

    Russell, "On Denoting"
    Russell, "Descriptions"
    Strawson, "On Referring"
    Donellan, "Reference and Definite Descriptions"
    Kripke, "Speaker's Reference and Semantic Reference"

	Segment 4 (weeks 7-8): Causal and description theories of reference

    Searle, "Proper Names"
    Kripke, Naming and Necessity, lectures I-II

	Segment 5 (weeks 9-10): Pragmatics

    Grice, "Logic and conversation"
    Stalnaker, "Pragmatic presuppositions"
